% ============================================================
%  Dispensa 02: Costruzione della Curva EIOPA Risk-Free Rate
%               Il metodo Smith-Wilson e la pipeline regolamentare
%  Corso di Calcolo Numerico — Laboratorio
%  Università degli Studi di Verona — A.A. 2026–2027
%  Riferimento: EIOPA-BoS-25-599 (Technical Documentation, Dec 2025)
% ============================================================
\documentclass[11pt,a4paper]{article}

\usepackage[T1]{fontenc}
\usepackage[utf8]{inputenc}
\usepackage[italian]{babel}
\usepackage{lmodern}
\usepackage{amsmath,amssymb,amsthm,mathtools}
\usepackage{geometry}
\usepackage{booktabs}
\usepackage{array}
\usepackage{enumitem}
\usepackage{graphicx}
\usepackage{xcolor,colortbl}
\usepackage{hyperref}
\usepackage{float}
\usepackage{fancyvrb}
\usepackage{tcolorbox}
\usepackage{titlesec}
\usepackage{bm}
\usepackage[ruled,vlined,linesnumbered]{algorithm2e}

\geometry{top=2.5cm,bottom=2.5cm,left=2.8cm,right=2.8cm}
\setlist[itemize]{topsep=4pt,itemsep=2pt,parsep=0pt}
\setlist[enumerate]{topsep=4pt,itemsep=2pt,parsep=0pt}
\hypersetup{colorlinks=true,linkcolor=blue!55!black,urlcolor=blue!55!black,
            citecolor=green!40!black}

\titleformat{\section}{\large\bfseries}{\thesection.}{0.55em}{}[{\vspace{1pt}\titlerule[0.4pt]}]
\titleformat{\subsection}{\normalsize\bfseries}{\thesubsection.}{0.55em}{}
\titleformat{\subsubsection}{\small\bfseries\itshape}{\thesubsubsection.}{0.5em}{}

% --- ambienti teorema ---
\theoremstyle{plain}
\newtheorem{teorema}{Teorema}[section]
\newtheorem{proposizione}[teorema]{Proposizione}
\theoremstyle{definition}
\newtheorem{definizione}[teorema]{Definizione}
\newtheorem{esempio}[teorema]{Esempio}
\theoremstyle{remark}
\newtheorem*{osservazione}{Osservazione}
\newtheorem*{attenzione}{Nota numerica}
\newtheorem{esercizio}{Esercizio}[section]

% --- box colorati ---
\tcbset{
  mybase/.style={
    colback=gray!5, colframe=gray!40,
    boxrule=0.4pt, arc=2pt,
    left=7pt, right=7pt, top=5pt, bottom=5pt,
    fonttitle=\small\bfseries,
    breakable
  }
}
\newtcolorbox{keybox}[1][Concetto chiave]{mybase,
  colframe=blue!40!black, colback=blue!4, title=#1}
\newtcolorbox{algobox}[1][Algoritmo]{mybase,
  colframe=teal!50!black, colback=teal!4, title=#1}
\newtcolorbox{warningbox}{mybase,
  colframe=orange!70!black, colback=orange!5}
\newtcolorbox{codebox}[1][Codice R]{mybase,
  colframe=green!40!black, colback=green!3, title=#1}

% --- colore righe tabella ---
\definecolor{rowtint}{RGB}{243,247,253}

% --- macro math ---
\newcommand{\R}{\mathbb{R}}
\newcommand{\UFR}{\mathrm{UFR}}
\newcommand{\LLP}{\mathrm{LLP}}
\newcommand{\CP}{\mathrm{CP}}
\newcommand{\CRA}{\mathrm{CRA}}
\newcommand{\bfzeta}{\bm{\zeta}}
\newcommand{\bfmu}{\bm{\mu}}
\newcommand{\Wmat}{\mathbf{W}}
\newcommand{\Hmat}{\mathbf{H}}
\newcommand{\abs}[1]{\left\lvert #1 \right\rvert}
\newcommand{\norm}[1]{\left\lVert #1 \right\rVert}
\newcommand{\defeq}{\coloneqq}

% ============================================================
\title{%
  \textbf{Costruzione della Curva EIOPA Risk-Free Rate}\\[6pt]
  \large Il Metodo Smith--Wilson: Sistema Lineare, Bisezione e\\[3pt]
  Interpolazione--Estrapolazione della Struttura a Termine\\[4pt]
  \normalsize Implementazione in R e Interpretazione Finanziaria
}
\author{Laboratorio di Calcolo Numerico --- Laurea Triennale in Matematica Applicata\\
  Universit\`a degli Studi di Verona\\[2pt]
  Anno Accademico 2026--2027}
\date{}

\begin{document}
\maketitle
\tableofcontents
\newpage

% ============================================================
\section{Introduzione e obiettivi}
\label{sec:intro}
% ============================================================

\subsection{Il problema}

Le compagnie di assicurazione vita assumono impegni che possono estendersi per
50--80 anni nel futuro: rendite vitalizie, polizze Long-Term Care, fondi pensione.
Per quantificarne il \emph{valore attuale}, occorre attualizzare i flussi di cassa
futuri stimati, il che richiede una \textbf{curva dei tassi di interesse} valida su
tutto l'orizzonte.

Il nodo cruciale \`e che i mercati finanziari forniscono tassi osservabili e affidabili
solo fino a una certa scadenza --- tipicamente 20 anni per l'euro. Oltre quel punto i
dati sono illiquidi o inesistenti. Il metodo di costruzione della curva deve quindi
risolvere due sotto-problemi simultaneamente:
\begin{enumerate}
  \item \textbf{Interpolazione}: ricostruire la curva tra i nodi di mercato osservati;
  \item \textbf{Estrapolazione}: prolungare la curva fino a 150 anni con convergenza
        a un tasso asintotico economicamente fondato (l'Ultimate Forward Rate).
\end{enumerate}

In questa dispensa mostriamo come l'intero problema si riduca a due passi numerici
ben definiti, che analizziamo in dettaglio:

\begin{enumerate}
  \item \textbf{Calibrazione --- sistema lineare denso $\Hmat\bfzeta = \mathbf{b}$
        (Sezione~\ref{sec:sistema}).}
        Smith--Wilson costruisce la curva di sconto come somma di un termine
        asintotico $e^{-\UFR_c t}$ e di $N$ funzioni ``kernel'' $W(t, u_j)$.
        I pesi $\zeta_j$ --- uno per ciascun nodo di mercato --- si ottengono
        imponendo che ogni swap par valga esattamente~1: questo genera un sistema
        lineare $N\times N$ denso, risolto con la fattorizzazione LU,
        costo $O(N^3)$.

  \item \textbf{Velocità di convergenza --- radice di $g(\alpha)=0$
        (Sezione~\ref{sec:alpha}).}
        Il kernel dipende da un parametro $\alpha > 0$ che controlla
        \emph{quanto velocemente} la curva converge verso l'UFR oltre il LLP.
        Un $\alpha$ piccolo $\Rightarrow$ convergenza lenta (la curva impiega
        molti anni prima di ``raggiungere'' l'UFR); un $\alpha$ grande
        $\Rightarrow$ convergenza rapida ma potenzialmente poco fedele ai dati.
        EIOPA fissa il criterio: $\alpha^*$ deve essere il \emph{minimo} valore
        tale che il tasso forward al Convergence Point ($\CP = 60$ anni) disti
        dall'UFR di meno di~1 bp, cioè
        \[
          g(\alpha) \defeq f(0, \CP;\, \alpha) - \UFR_c = 0\;\text{(a 1 bp)}.
        \]
        Trovare $\alpha^*$ equivale quindi a trovare la \textbf{radice} di
        $g$, che dipende non linearmente da $\alpha$ (per ogni candidato
        occorre risolvere nuovamente il sistema lineare). Si usa il metodo
        di bisezione per la sua convergenza garantita.
\end{enumerate}

\begin{keybox}[Il filo rosso: da UFR a numeri a curva]
\textbf{Passo 0 (dato economico)} Si parte dall'UFR --- il tasso forward di
lungo periodo stima dai mercati ($3.30\%$ per EUR 2025).

\textbf{Passo 1 (algebra lineare)} Dati i $N=12$ tassi OIS di mercato
(corretto per il CRA), si risolve $\Hmat(\alpha)\bfzeta = \mathbf{1} - \bfmu_Q$:
si ottengono i pesi $\zeta_j$ che ``incollano'' la curva ai dati.

\textbf{Passo 2 (ricerca di zeri)} Si cerca $\alpha^*$ tale che la curva così
calibrata converga all'UFR con la velocità prescritta da EIOPA. Ogni
valutazione di $g(\alpha)$ richiede di rifare il Passo~1, quindi i due
problemi sono intrinsecamente accoppiati.

\textbf{Risultato} Un'unica formula analitica $P(t) = e^{-\UFR_c t} +
\sum_j \zeta_j W(t,u_j)$ valida per $t \in [0,150]$, che interpola i dati,
converge all'UFR e rispetta tutti i vincoli normativi.
\end{keybox}

\subsection{Finalit\`a didattiche}

Questa dispensa illustra un caso reale in cui i metodi del calcolo numerico
svolgono un ruolo diretto nel settore finanziario--assicurativo:

\begin{center}
\renewcommand{\arraystretch}{1.3}
\begin{tabular}{lll}
\toprule
\textbf{Metodo numerico} & \textbf{Applicazione EIOPA} & \textbf{Riferimento}\\
\midrule
Sistema lineare denso (LU) & Calibrazione Smith--Wilson & sez.~9.15\\
Bisezione / ricerca di zeri & $\alpha^*$: velocità conv.\ curva $\to$ UFR & sez.~9.4\\
Numero di condizionamento & Stabilit\`a della curva & analisi degli errori\\
\emph{Confronto:} Thomas $O(N)$ & Spline cubica (alternativa) & de Boor~\cite{deBoor1978}\\
\bottomrule
\end{tabular}
\end{center}

Il codice R di supporto \`e nel file \texttt{R/curva\_eiopa\_rfr\_CN.R} e produce
tutti i grafici mostrati nella dispensa.

\subsection{Il quadro normativo: Solvency~II e EIOPA}

Il quadro normativo europeo \textbf{Solvency~II} (Direttiva 2009/138/CE~\cite{dir2009},
art.~77e) richiede alle imprese assicurative di calcolare le riserve tecniche
attualizzando con una curva risk-free ufficiale. L'\textbf{EIOPA}
(European Insurance and Occupational Pensions Authority) pubblica questa struttura
a termine \emph{mensilmente} per tutte le valute rilevanti.

I principi guida (sez.~1--3 di \cite{eiopa2025}):
\begin{itemize}
  \item \textbf{Replicabilit\`a}: ogni impresa deve poter riprodurre la curva;
  \item \textbf{Coerenza di mercato}: si usano solo dati da mercati DLT
        (deep, liquid, transparent);
  \item \textbf{Stabilit\`a}: la variazione dei parametri chiave \`e limitata.
\end{itemize}

% ------------------------------------------------------------
\subsection{Struttura dei tassi di interesse: spot, forward e par}
\label{sec:tassi-def}
% ------------------------------------------------------------

Per comprendere l'input e l'output del metodo Smith--Wilson occorre distinguere
tre rappresentazioni diverse della struttura a termine, tutte derivabili dalla
stessa funzione di sconto $P(t)$ ma utili in contesti diversi.

\begin{definizione}[Funzione di sconto e tasso spot]
La \textbf{funzione di sconto} $P(t)>0$ \`e il prezzo odierno di $1$~euro
pagato con certezza all'istante $t \ge 0$.
Il corrispondente \textbf{tasso spot} a capitalizzazione continua \`e:
\[
  r(0,t) \defeq -\frac{\ln P(t)}{t}, \qquad t > 0,
\]
ovvero $P(t) = e^{-r(0,t)\,t}$.
La struttura a termine $t \mapsto r(0,t)$ per $t=1,\ldots,150$ anni \`e
esattamente ci\`o che EIOPA pubblica nel foglio \texttt{RFR\_spot\_no\_VA}.
\end{definizione}

\begin{definizione}[Tasso forward istantaneo]
Il \textbf{tasso forward istantaneo} alla scadenza $t$ \`e il rendimento
marginale dell'investimento nell'intorno di $t$:
\[
  f(0,t) \defeq -\frac{\mathrm{d}}{\mathrm{d}t}\ln P(t)
               = r(0,t) + t\,r'(0,t).
\]
Economicamente, $f(0,t)\,\mathrm{d}t$ \`e il rendimento di un impegno
istantaneo concordato oggi per la finestra $[t,t+\mathrm{d}t]$.
Il criterio di convergenza EIOPA \`e enunciato proprio in questi termini:
si richiede che il tasso forward al punto di convergenza
($\CP=60$ anni) disti dall'UFR di meno di 1~bp, cio\`e
$|f(0,\CP;\,\alpha^*) - \UFR_c| < 0.0001$.
\end{definizione}

\begin{definizione}[Tasso par di uno swap]
\label{def:tasso_par}
Uno \textbf{swap su tassi} (interest rate swap) \`e un contratto in cui due
controparti si scambiano periodicamente interessi calcolati sullo stesso
nozionale (di norma non scambiato):
\begin{itemize}
  \item \textbf{gamba fissa}: paga un tasso costante $s$;
  \item \textbf{gamba variabile}: paga un tasso indicizzato (qui overnight \texteuro{}STR composto).
\end{itemize}
Alla data iniziale si scambia solo il \emph{valore netto} delle due gambe.

Siano $T_1 < T_2 < \cdots < T_n$ le date di pagamento di uno swap
(tipicamente annuali, $\Delta_k \approx 1$).
Qui $\Delta_k$ indica la \textbf{frazione d'anno} (year fraction) tra
$T_{k-1}$ e $T_k$ secondo la day-count convention del contratto:
\[
  \Delta_k \defeq \text{YearFrac}(T_{k-1},T_k)
  \approx T_k-T_{k-1}
\]
(in anni, se il tempo \`e gi\`a espresso in unit\`a annuali).
Negli OIS annuali usati qui, tipicamente $\Delta_k \simeq 1$ per ogni $k$
(cfr. Hull~\cite{hull2018}, cap.~7 e cap.~28).

Il \textbf{tasso par} $s_n$ alla scadenza $T_n$ \`e il tasso fisso che
rende nullo il valore netto iniziale del contratto:
\begin{equation}
  \label{eq:tasso_par}
  s_n \cdot \sum_{k=1}^{n} \Delta_k\,P(T_k) + P(T_n) = 1
  \quad\Longrightarrow\quad
  s_n = \frac{1 - P(T_n)}{\displaystyle\sum_{k=1}^{n} \Delta_k\,P(T_k)}.
\end{equation}
Il numeratore $1-P(T_n)$ \`e il guadagno netto sulla gamba variabile;
il denominatore \`e il valore attuale di una rendita unitaria.
Il nome ``par'' deriva dal fatto che un titolo con cedola $s_n$ e rimborso
del nozionale 1 alla scadenza $T_n$ vale esattamente 1 (``alla pari'') oggi.
\end{definizione}

\begin{keybox}[Perch\'e Smith--Wilson usa i tassi par e non quelli spot?]
I tassi spot $r(0,t)$ \emph{non} sono osservabili direttamente sul mercato.
Il metodo tradizionale per ricavarli \`e il \textbf{bootstrap}: risolvere
la~\eqref{eq:tasso_par} sequenzialmente per $n=1,2,\ldots,N$, estraendo
$P(T_n)$ a ogni passo. Il bootstrap introduce per\`o propagazione degli
errori e richiede un'interpolazione aggiuntiva tra i nodi.

Smith--Wilson elimina questo passaggio intermedio: usa \emph{direttamente}
i tassi par osservati come vincoli di calibrazione e ricostruisce $P(t)$
(e quindi $r(0,t)$ e $f(0,t)$) in forma chiusa tramite la
formula~\eqref{eq:sw_main}.
Questo \`e quanto afferma esplicitamente la sez.~9.1 di~\cite{eiopa2025}:
\emph{``The Smith-Wilson method [\ldots] is applied directly to par swap rates''}.
\end{keybox}

\subsection{Origine dei dati}
\label{sec:dati-origine}

I dati di input per la curva EIOPA EUR provengono da due fonti:

\paragraph{(a) Tassi OIS su \texteuro{}STR.}
Gli \textbf{Overnight Index Swap (OIS)} riferiti a \texteuro{}STR
(Euro Short-Term Rate, pubblicato da BCE alle 8:00 CET)
sono il proxy standard del tasso privo di rischio per l'area euro.
I provider di dati utilizzati da EIOPA sono \textbf{Bloomberg} (ticker: EESWE*)
e \textbf{Refinitiv/LSEG}, come specificato nell'Annex~B di~\cite{eiopa2025}.

Un OIS a scadenza $T$ \`e un contratto derivato in cui:
\begin{itemize}
  \item una parte paga il \textbf{tasso fisso} $s$ (tasso par);
  \item l'altra parte paga il tasso overnight composto per la durata del contratto.
\end{itemize}
Non c'\`e scambio di capitale: il rischio di credito \`e quasi nullo
(solo sul mark-to-market del derivato, non sul nozionale).

\paragraph{(b) Pubblicazione mensile EIOPA.}
EIOPA pubblica il file \emph{EIOPA\_RFR\_yyyymmdd.xlsx} il primo giorno lavorativo
di ogni mese, scaricabile da
\url{https://www.eiopa.europa.eu/tools-and-data/risk-free-interest-rate-term-structures}.
Il file contiene:
\begin{itemize}
  \item foglio \texttt{RFR\_spot\_no\_VA}: curva spot senza Volatility Adjustment;
  \item foglio \texttt{Spot\_NO\_VA\_shock\_UP/DOWN}: curve per SCR (stress);
  \item parametri utilizzati ($\alpha^*$, UFR, LLP, CRA) nel foglio \texttt{Parameters}.
\end{itemize}

\paragraph{(c) Dati utilizzati in questa dispensa.}
Per ragioni didattiche utilizziamo un set di 12 tassi OIS par illustrativi
ma rappresentativi del livello di mercato EUR a maggio 2025. Le scadenze
standard DLT per l'EUR sono:
\[
  T \in \{1, 2, 3, 4, 5, 6, 7, 8, 10, 12, 15, 20\} \text{ anni}.
\]
Le scadenze 9, 11, 13, 14, 16--19 non sono incluse perch\'e non soddisfano
i criteri di liquidit\`a (residual volume $<6\%$ del totale, sez.~3.7 EIOPA).

\subsection{La pipeline EIOPA in cinque passi}

\begin{center}\small
\renewcommand{\arraystretch}{1.3}
\begin{tabular}{clp{7.5cm}}
\toprule
\textbf{Fase} & \textbf{Nome} & \textbf{Operazione} \\
\midrule
1 & Raccolta dati & Tassi OIS (\texteuro{}STR) da Bloomberg/Refinitiv \\
2 & Valutazione DLT & Verifica liquidit\`a; determinazione del \textbf{LLP} \\
3 & CRA & Sottrazione del Credit Risk Adjustment (10 bps per EUR 2025) \\
4 & \textbf{Calibrazione SW} & Smith--Wilson applicato direttamente ai swap par: \\
  & & risoluzione del sistema lineare $\Hmat\bfzeta = \mathbf{1}-\bfmu_Q$ \\
5 & Parametro $\alpha$ & Ricerca della radice di $g(\alpha) = f(0,\CP;\alpha) - \UFR_c$ \\
\bottomrule
\end{tabular}
\end{center}

Il nucleo computazionale \`e interamente nelle fasi 4 e 5.

\begin{warningbox}
\textbf{Cosa NON fa EIOPA.}
La sez.~9.1 di~\cite{eiopa2025} afferma esplicitamente:
\emph{``The interpolation, where necessary, and extrapolation of interest rates
have been developed applying the Smith-Wilson method.''}
Non c'\`e alcun passo intermedio di bootstrap n\'e di spline cubiche nella
pipeline ufficiale. Smith--Wilson viene applicato \emph{direttamente} ai tassi par
degli swap (sez.~9.15).
\end{warningbox}

% ------------------------------------------------------------
\subsection{Dati reali: fonti e procedura per ricostruire la curva EUR maggio 2026}
\label{sec:dati_reali}
% ------------------------------------------------------------

Per ricostruire \emph{esattamente} la curva EIOPA EUR di un dato mese occorrono
tre categorie di input. Di seguito indichiamo le fonti ufficiali e le alternative
accessibili per il mese di \textbf{maggio 2026}.

\begin{keybox}[Date di riferimento --- maggio 2026]
\begin{itemize}
  \item \textbf{Data di riferimento (valuation date)}: venerd\`i \textbf{29 maggio 2026}
        (ultimo giorno lavorativo del mese; il 30~maggio \`e sabato).
  \item \textbf{Data di pubblicazione}: luned\`i \textbf{1\textordmasculine{} giugno 2026}
        (primo giorno lavorativo del mese successivo).
  \item \textbf{Nome file atteso}: \texttt{EIOPA\_RFR\_20260529.xlsx}
\end{itemize}
\end{keybox}

% ......................... FONTE A .........................
\subsubsection*{Fonte~A --- File ufficiale EIOPA (per la verifica del risultato)}

\noindent\textbf{URL:}
\url{https://www.eiopa.europa.eu/tools-and-data/risk-free-interest-rate-term-structures\_en}

\noindent Scaricare il file \texttt{EIOPA\_RFR\_20260529.xlsx}.
I fogli rilevanti sono i seguenti:

\begin{center}\small
\renewcommand{\arraystretch}{1.25}
\begin{tabular}{lp{9cm}}
\toprule
\textbf{Foglio Excel} & \textbf{Contenuto} \\
\midrule
\texttt{RFR\_spot\_no\_VA}  & Tassi spot ufficiali per $T=1,\ldots,150$ anni (senza
                               Volatility Adjustment). \textbf{Usare come confronto
                               finale.} \\
\texttt{Parameters}         & $\alpha^*$, UFR, LLP, CRA effettivamente usati da EIOPA
                               per questo mese. \\
\texttt{Spot\_NO\_VA\_shock\_UP} \newline
\texttt{Spot\_NO\_VA\_shock\_DOWN} & Curve sotto stress istantaneo per il calcolo
                               del Solvency Capital Requirement (SCR). \\
\bottomrule
\end{tabular}
\end{center}

\begin{warningbox}
Il file EIOPA contiene la curva \emph{risultato}, non i tassi OIS di input.
Serve per \textbf{verificare} la ricostruzione ma non pu\`o sostituire i dati
di mercato (Fonte~B) se si vuole replicare il calcolo dall'inizio.
\end{warningbox}

% ......................... FONTE B .........................
\subsubsection*{Fonte~B --- Tassi OIS \texteuro{}STR par: 12 scadenze DLT}

Questi sono i dati di mercato veri e propri, rilevati alla chiusura del
29~maggio~2026. Occorrono i tassi par degli OIS (Overnight Index Swap)
referenziati a \texteuro{}STR per le 12 scadenze DLT:
$\{1,2,3,4,5,6,7,8,10,12,15,20\}$ anni.

\paragraph{B1 --- Bloomberg Terminal (fonte primaria EIOPA, Annex~B di~\cite{eiopa2025}).}
I ticker da interrogare con il comando \texttt{HP <GO>} (Historical Price)
alla data 29/05/2026, campo \texttt{PX\_LAST}:

\begin{center}\small
\renewcommand{\arraystretch}{1.2}
\begin{tabular}{clcl}
\toprule
\textbf{Scadenza} & \textbf{Ticker Bloomberg} & \textbf{Scadenza} & \textbf{Ticker Bloomberg} \\
\midrule
1Y  & \texttt{EUSWEC1 BG Curncy}  & 8Y  & \texttt{EUSWEC8 BG Curncy}  \\
2Y  & \texttt{EUSWEC2 BG Curncy}  & 10Y & \texttt{EUSWEC10 BG Curncy} \\
3Y  & \texttt{EUSWEC3 BG Curncy}  & 12Y & \texttt{EUSWEC12 BG Curncy} \\
4Y  & \texttt{EUSWEC4 BG Curncy}  & 15Y & \texttt{EUSWEC15 BG Curncy} \\
5Y  & \texttt{EUSWEC5 BG Curncy}  & 20Y & \texttt{EUSWEC20 BG Curncy} \\
6Y  & \texttt{EUSWEC6 BG Curncy}  &     &                              \\
7Y  & \texttt{EUSWEC7 BG Curncy}  &     &                              \\
\bottomrule
\end{tabular}
\end{center}

\noindent Il prefisso \texttt{EUSWEC} sta per ``EUR Swap \texteuro{}STR Compounded''.
La serie storica \`e disponibile dall'ottobre~2019 (avvio di \texteuro{}STR).
Per il download massivo: \texttt{DAPI <GO>} $\to$ Excel Add-in, oppure
direttamente \texttt{=BDH("EUSWEC1 BG Curncy","PX\_LAST","29/05/2026","29/05/2026")}.

\paragraph{B2 --- Refinitiv Workspace / LSEG Eikon (fonte alternativa EIOPA).}
I RIC (Reuters Instrument Code) corrispondenti sono:

\begin{center}\small
\renewcommand{\arraystretch}{1.2}
\begin{tabular}{clcl}
\toprule
\textbf{Scadenza} & \textbf{RIC Refinitiv} & \textbf{Scadenza} & \textbf{RIC Refinitiv} \\
\midrule
1Y  & \texttt{EURESTOIS1Y=}  & 8Y  & \texttt{EURESTOIS8Y=}  \\
2Y  & \texttt{EURESTOIS2Y=}  & 10Y & \texttt{EURESTOIS10Y=} \\
3Y  & \texttt{EURESTOIS3Y=}  & 12Y & \texttt{EURESTOIS12Y=} \\
4Y  & \texttt{EURESTOIS4Y=}  & 15Y & \texttt{EURESTOIS15Y=} \\
5Y  & \texttt{EURESTOIS5Y=}  & 20Y & \texttt{EURESTOIS20Y=} \\
6Y  & \texttt{EURESTOIS6Y=}  &     &                        \\
7Y  & \texttt{EURESTOIS7Y=}  &     &                        \\
\bottomrule
\end{tabular}
\end{center}

\noindent In Eikon/Workspace: funzione \texttt{TR()} in Excel con campo \texttt{TR.FIXEDRATE}
e parametro \texttt{SDate=2026-05-29}.

\paragraph{B3 --- ECB Data Portal (fonte libera, parziale).}
La BCE pubblica il tasso overnight \texteuro{}STR ma \emph{non} i tassi OIS
par a termine. \`E tuttavia possibile accedere ai dati del \textbf{mercato
monetario} (dataset \texttt{IVF}) e ai dati storici dell'OIS dalla pagina:
\url{https://data.ecb.europa.eu} $\to$ ``Money Market Statistics''.
Si tratta di dati complementari; per i tassi OIS par a termine (scadenze 1--20Y)
rimane necessario Bloomberg o Refinitiv.

\begin{attenzione}[Nota sul mercato OIS \texteuro{}STR]
Prima di ottobre~2019 i tassi erano referenziati a EONIA. Da quella data
\texteuro{}STR (\textit{Euro Short-Term Rate}) \`e il tasso overnight di
riferimento BCE (art.~5 del Reg.~2019/1011). I ticker Bloomberg \texttt{EUSWEC*}
coprono la transizione; per dati pre-2019 usare i ticker \texttt{EUSWE*}
(EONIA-based) e verificare la sovrapposizione nel periodo di transizione.
\end{attenzione}

% ......................... FONTE C .........................
\subsubsection*{Fonte~C --- Parametri EIOPA (UFR, LLP, CRA, $\alpha^*$)}

\paragraph{C1 --- UFR per EUR 2026.}
L'UFR viene rivisto annualmente da EIOPA (processo di cui alla sez.~9.3
di~\cite{eiopa2025}). Il valore per il 2026 \`e pubblicato nel documento:
\begin{itemize}
  \item \textbf{Documento}: \emph{EIOPA Decision on the UFR for 2026}
  \item \textbf{URL}: \url{https://www.eiopa.europa.eu/publications} $\to$
        ricerca ``UFR 2026''
\end{itemize}
Il valore per EUR, storicamente stabile a $3.30\%$ dal 2023, va verificato
ogni anno prima di procedere al calcolo.

\paragraph{C2 --- LLP e CRA per EUR.}
Pubblicati mensilmente nel \textbf{Technical Information} della BCE e nel
foglio \texttt{Parameters} del file EIOPA:
\begin{itemize}
  \item LLP EUR: \textbf{20 anni} (stabile; dipende dall'analisi DLT semestrale)
  \item CRA EUR: da verificare nel foglio \texttt{Parameters} ---
        attualmente $\mathbf{10}$ bps (Situazione~1, sez.~7.3 di~\cite{eiopa2025})
\end{itemize}

\paragraph{C3 --- $\alpha^*$ calcolato vs pubblicato.}
Il valore di $\alpha^*$ non viene fornito come input: si calcola con la bisezione
(Sezione~\ref{sec:alpha}). Tuttavia il foglio \texttt{Parameters} del file EIOPA
riporta il valore usato ufficialmente, utile come \textbf{verifica}.

% ......................... TABELLA RACCOLTA .........................
\subsubsection*{Tabella riepilogativa: dati da raccogliere per maggio 2026}

\begin{center}\small
\renewcommand{\arraystretch}{1.3}
\begin{tabular}{clp{4cm}l}
\toprule
\textbf{Input} & \textbf{Dato} & \textbf{Fonte} & \textbf{Valore da rilevare} \\
\midrule
$s_1$  & OIS 1Y  par & Bloomberg \texttt{EUSWEC1}  & \underline{\hspace{2cm}}\% \\
$s_2$  & OIS 2Y  par & Bloomberg \texttt{EUSWEC2}  & \underline{\hspace{2cm}}\% \\
$s_3$  & OIS 3Y  par & Bloomberg \texttt{EUSWEC3}  & \underline{\hspace{2cm}}\% \\
$s_4$  & OIS 4Y  par & Bloomberg \texttt{EUSWEC4}  & \underline{\hspace{2cm}}\% \\
$s_5$  & OIS 5Y  par & Bloomberg \texttt{EUSWEC5}  & \underline{\hspace{2cm}}\% \\
$s_6$  & OIS 6Y  par & Bloomberg \texttt{EUSWEC6}  & \underline{\hspace{2cm}}\% \\
$s_7$  & OIS 7Y  par & Bloomberg \texttt{EUSWEC7}  & \underline{\hspace{2cm}}\% \\
$s_8$  & OIS 8Y  par & Bloomberg \texttt{EUSWEC8}  & \underline{\hspace{2cm}}\% \\
$s_{10}$ & OIS 10Y par & Bloomberg \texttt{EUSWEC10} & \underline{\hspace{2cm}}\% \\
$s_{12}$ & OIS 12Y par & Bloomberg \texttt{EUSWEC12} & \underline{\hspace{2cm}}\% \\
$s_{15}$ & OIS 15Y par & Bloomberg \texttt{EUSWEC15} & \underline{\hspace{2cm}}\% \\
$s_{20}$ & OIS 20Y par & Bloomberg \texttt{EUSWEC20} & \underline{\hspace{2cm}}\% \\
\midrule
$\UFR_{\rm ann}$ & UFR EUR 2026 & EIOPA Decision UFR 2026  & \underline{\hspace{2cm}}\% \\
$\CRA$ & Credit Risk Adj. & EIOPA file, foglio \texttt{Parameters} & \underline{\hspace{2cm}} bps \\
$\alpha^*_{\rm EIOPA}$ & Parametro $\alpha$ ufficiale & EIOPA file, foglio \texttt{Parameters} & \underline{\hspace{2cm}} \\
\bottomrule
\end{tabular}
\end{center}

% ......................... VERIFICA .........................
\subsubsection*{Procedura di verifica in R}

Una volta raccolti i dati, la pipeline di verifica in R \`e:

\begin{enumerate}
  \item Inserire i 12 tassi $s_k$ (lordi) nel vettore \texttt{s\_mkt}.
  \item Applicare il CRA: \texttt{r\_mkt <- s\_mkt - CRA}.
  \item Eseguire \texttt{R/curva\_eiopa\_rfr\_CN.R} (Sezione~\ref{sec:codice}).
  \item Confrontare \texttt{alpha\_star} calcolato con $\alpha^*_{\rm EIOPA}$ dal
        foglio \texttt{Parameters}: la differenza deve essere $< 10^{-6}$.
  \item Confrontare la curva spot $r(0,T)$ per $T=1,\ldots,150$ con il foglio
        \texttt{RFR\_spot\_no\_VA}: l'errore massimo deve essere $< 0.1$ bps
        (vedi Esercizio~5 per il codice di confronto).
\end{enumerate}

% ============================================================
\section{Prerequisiti matematici}
\label{sec:prerequisiti}
% ============================================================

In questa sezione presentiamo gli strumenti matematici necessari per comprendere
l'algoritmo Smith--Wilson, organizzati attorno a due problemi fondamentali:

\begin{keybox}[I due problemi fondamentali]
\begin{enumerate}
  \item \textbf{Sistema lineare}: data $\Hmat \in \R^{N\times N}$ densa, trovare
        $\bfzeta$ tale che $\Hmat\bfzeta = \mathbf{b}$.
  \item \textbf{Ricerca di zeri}: trovare $\alpha^* \in [a,b]$ tale che
        $g(\alpha^*) \approx 0$.
\end{enumerate}
\end{keybox}

La risposta al primo problema passa attraverso la fattorizzazione LU; la risposta
al secondo attraverso il metodo di bisezione. Entrambi i metodi forniscono
\emph{garanzie} sulla soluzione (esistenza, unicità, stabilità).

% ------------------------------------------------------------
\subsection{Sistemi lineari e fattorizzazione LU}
\label{sec:lu}
% ------------------------------------------------------------

\begin{definizione}[Fattorizzazione LU]
Data $A \in \R^{n\times n}$, una fattorizzazione LU consiste nel trovare una
matrice triangolare inferiore $L$ (con $l_{ii}=1$) e una triangolare superiore $U$
tali che $A = LU$. In generale, per garantire esistenza e stabilit\`a numerica si
introduce una matrice di permutazione $P_\pi$:
\[
  P_\pi A = LU.
\]
\end{definizione}

\begin{teorema}[Esistenza con pivoting parziale]
\label{thm:lu}
Per ogni matrice $A \in \R^{n\times n}$ non singolare, esiste una matrice di
permutazione $P_\pi$ tale che $P_\pi A = LU$ con $L$ triangolare inferiore unitaria
e $U$ triangolare superiore.
\end{teorema}

\paragraph{Algoritmo.}
La fattorizzazione si ottiene con l'\textbf{eliminazione gaussiana con pivoting parziale}:
ad ogni passo $k=1,\ldots,n-1$, si scambia la riga con il pivot di modulo massimo nella
colonna $k$ (dalla riga $k$ in gi\`u), poi si eliminano gli elementi sotto-diagonali.

\paragraph{Costo computazionale.}
Il numero di operazioni floating-point per la fattorizzazione \`e:
\begin{equation}
  \label{eq:costo_lu}
  C_{\text{LU}} = \frac{2}{3}n^3 + O(n^2).
\end{equation}
La risoluzione dei due sistemi triangolari ($Ly=Pb$, $Ux=y$) costa $O(n^2)$ ciascuno.
Per $n=12$ (numero di nodi EIOPA EUR): $C \approx 1\,152$ operazioni.

\begin{proposizione}[Stabilit\`a della fattorizzazione LU con pivoting]
\label{prop:stab_lu}
Sia $A\in\R^{n\times n}$ e sia $\hat{x}$ la soluzione calcolata in aritmetica
floating-point con unit\`a di arrotondamento $u$. Allora:
\[
  \frac{\norm{x - \hat{x}}}{\norm{x}} \le c(n)\,\kappa(A)\,u,
\]
dove $\kappa(A) = \norm{A}\norm{A^{-1}}$ \`e il numero di condizionamento e
$c(n)$ \`e una costante moderata (limitata dal fattore di crescita $\rho_n$,
tipicamente $\rho_n \ll 2^{n-1}$ in pratica).
\end{proposizione}

% ------------------------------------------------------------
\subsection{Scelta del metodo risolutivo: perch\`e LU?}
\label{sec:scelta_metodo}
% ------------------------------------------------------------

Prima di procedere \`e utile motivare la scelta della fattorizzazione LU
rispetto alle alternative classiche del calcolo numerico.

\begin{center}\small
\renewcommand{\arraystretch}{1.25}
\begin{tabular}{p{2.4cm}p{3.2cm}p{5.8cm}c}
\toprule
\textbf{Metodo} & \textbf{Requisiti} & \textbf{Limitazione per Smith--Wilson} & \textbf{OK?}\\
\midrule
\textbf{LU con pivoting} &
  $A$ quadrata non singolare &
  Nessuna &
  \checkmark\\
\addlinespace
Cholesky &
  $A$ simmetrica def.\ pos. &
  $\Hmat$ \`e densa ma \textbf{non simmetrica}: righe $=$ swap per $T_i$, colonne $=$ nodi kernel $u_k$ &
  $\times$\\
\addlinespace
QR (Householder) &
  Qualsiasi $A$ &
  Costo $\approx 2n^3$, il doppio di LU; la stabilit\`a aggiuntiva \`e superflua per $\kappa\approx 100$ &
  $\sim$\\
\addlinespace
Metodi iterativi\newline(Jacobi, CG, GMRES) &
  Matrici sparse o molto grandi &
  Per $N=12$ il costo di setup supera il tempo di risoluzione; la convergenza richiede precondizionatori non banali &
  $\times$\\
\bottomrule
\end{tabular}
\end{center}

\begin{keybox}[Regola pratica: LU per sistemi densi piccoli--medi]
Per sistemi densi di dimensioni $N \lesssim 500$, LU con pivoting parziale \`e
lo standard industriale: \`e la prima scelta di MATLAB (\texttt{\textbackslash}),
R (\texttt{solve()}) e NumPy (\texttt{linalg.solve()}), tutti basati su
LAPACK \texttt{dgesv}. I metodi iterativi diventano convenienti solo quando
$N \gtrsim 10^3$ \emph{e} la matrice \`e \textbf{sparsa}. Smith--Wilson con
$N=12$ \`e un caso emblematico: la robustezza e la semplicita\`a di LU
sono pi\`u importanti dell'efficienza asintotica.
\end{keybox}

% ------------------------------------------------------------
\subsection{Numero di condizionamento}
\label{sec:condizionamento}
% ------------------------------------------------------------

\begin{definizione}[Numero di condizionamento]
Il numero di condizionamento in norma 2 di una matrice $A$ non singolare \`e:
\[
  \kappa_2(A) = \norm{A}_2 \cdot \norm{A^{-1}}_2 = \frac{\sigma_{\max}(A)}{\sigma_{\min}(A)}.
\]
\end{definizione}

Questo numero controlla l'amplificazione degli errori: una perturbazione relativa
$\varepsilon$ nei dati pu\`o causare una perturbazione relativa fino a
$\kappa_2(A)\cdot\varepsilon$ nella soluzione.

\begin{keybox}[Condizionamento nel contesto EIOPA]
Per la matrice $\Hmat$ di Smith--Wilson con 12 nodi EUR e $\alpha \approx 0.128$,
si ha tipicamente $\kappa_2(\Hmat) \approx 10$--$100$: il sistema \`e
\textbf{ben condizionato}. Ma per $\alpha \to 0.05$ (vicino al limite inferiore),
le colonne di $\Hmat$ diventano quasi proporzionali e $\kappa_2$ pu\`o
crescere di 1--2 ordini di grandezza (Figura~\ref{fig:cond}).

\smallskip
\textbf{Interpretazione pratica.}
Con $\kappa_2(\Hmat) \approx 100$ e unit\`a di arrotondamento IEEE~754
$u \approx 2.2\times 10^{-16}$ (doppia precisione), la
Proposizione~\ref{prop:stab_lu} garantisce:
\[
  \frac{\norm{\bfzeta - \hat{\bfzeta}}}{\norm{\bfzeta}}
  \;\lesssim\; \kappa_2(\Hmat)\cdot u
  \;\approx\; 100 \times 10^{-16} = 10^{-14}.
\]
Questo \`e \textbf{10 ordini di grandezza} al di sotto della precisione
richiesta da EIOPA ($1\,\mathrm{bp} = 10^{-4}$): il sistema lineare \`e
risolto in modo praticamente esatto dalla fattorizzazione LU a doppia precisione.
\end{keybox}

% ------------------------------------------------------------
\subsection{Il metodo di bisezione}
\label{sec:bisezione}
% ------------------------------------------------------------

\begin{teorema}[Esistenza di zeri --- Bolzano]
Sia $g: [a,b] \to \R$ continua con $g(a)\cdot g(b) < 0$. Allora esiste
$\alpha^* \in (a,b)$ tale che $g(\alpha^*) = 0$.
\end{teorema}

\begin{definizione}[Metodo di bisezione]
Dato un intervallo $[a,b]$ con $g(a)\cdot g(b) < 0$, il metodo di bisezione
genera una successione di intervalli $[a_n, b_n]$ con:
\[
  m_n = \frac{a_n + b_n}{2}, \qquad
  [a_{n+1}, b_{n+1}] = \begin{cases}
    [a_n, m_n] & \text{se } g(a_n)\cdot g(m_n) < 0,\\
    [m_n, b_n] & \text{altrimenti.}
  \end{cases}
\]
\end{definizione}

\begin{proposizione}[Convergenza della bisezione]
\label{prop:bisez}
Dopo $n$ iterazioni su un intervallo iniziale di ampiezza $\ell = b-a$:
\begin{equation}
  \label{eq:bisez_err}
  \abs{\alpha^* - m_n} \le \frac{\ell}{2^{n+1}}.
\end{equation}
La convergenza \`e \textbf{lineare} con rapporto $1/2$. Per raggiungere una
tolleranza $\varepsilon$ servono:
\begin{equation}
  \label{eq:bisez_iter}
  n \ge \left\lceil \log_2\!\left(\frac{\ell}{\varepsilon}\right) \right\rceil
  \text{ iterazioni.}
\end{equation}
\end{proposizione}

\begin{attenzione}[Bisezione vs metodi pi\`u rapidi]
La bisezione ha convergenza solo lineare ($O(1/2^n)$), ma \`e:
\begin{itemize}
  \item \textbf{globalmente convergente} (non richiede una buona stima iniziale);
  \item \textbf{priva di derivate} (non serve $g'$);
  \item \textbf{con costo prevedibile} (il numero di iterazioni \`e noto a priori).
\end{itemize}
Per il parametro $\alpha$ EIOPA queste propriet\`a sono decisive: la funzione
$g(\alpha) = f(0,\CP;\alpha) - \UFR_c$ non ha derivata analitica semplice (dipende
dalla soluzione di un sistema lineare) e 20 iterazioni sono pi\`u che sufficienti.
\end{attenzione}

\begin{keybox}[Perch\`e non Newton--Raphson?]
Newton converge quadraticamente ($\sim c^{2^n}$, $c<1$) e potrebbe sembrare
più rapido della bisezione. Per $g(\alpha)$ tuttavia la derivata $g'(\alpha)$
non \`e disponibile analiticamente: differenziare $g$ rispetto ad $\alpha$
equivale a differenziare la soluzione $\bfzeta(\alpha)$ del sistema lineare,
il che richiede un ulteriore sistema di dimensione $N$.

\noindent Con derivata \textbf{numerica centrata} ($h=10^{-4}$):
\begin{center}\small
\renewcommand{\arraystretch}{1.2}
\begin{tabular}{lccc}
\toprule
\textbf{Metodo} & \textbf{LU/iter.} & \textbf{Iterazioni} & \textbf{LU totali}\\
\midrule
Bisezione & 1 & $\approx 20$ & $\mathbf{20}$\\
Newton (deriv.\ numerica) & 3\;(2 per $g'$, 1 per $g$) & $\approx 7$ & $\approx 21$\\
\bottomrule
\end{tabular}
\end{center}
I costi totali sono quasi identici, ma la bisezione offre convergenza
\emph{garantita} senza bisogno di una stima iniziale vicina a $\alpha^*$.
Per questo tipo di problema --- funzione liscia, intervallo noto, nessun costo
aggiuntivo per iterazione --- la bisezione \`e la scelta ingegneristicamente
corretta.
\end{keybox}

% ------------------------------------------------------------
\subsection{Confronto: sistema tridiagonale e algoritmo di Thomas}
\label{sec:thomas}
% ------------------------------------------------------------

Per le spline cubiche (presentate come confronto nella Sezione~\ref{sec:confronto})
il sistema lineare ha una struttura speciale: la matrice \`e \textbf{tridiagonale}.

\begin{definizione}[Matrice tridiagonale]
$A \in \R^{n\times n}$ \`e tridiagonale se $a_{ij} = 0$ per $|i-j|>1$.
\end{definizione}

\begin{proposizione}[Algoritmo di Thomas --- costo $O(n)$]
\label{prop:thomas}
Se $A$ \`e tridiagonale e a dominanza diagonale stretta, la fattorizzazione LU
(senza pivoting) esiste, \`e stabile e costa esattamente $O(n)$ operazioni
(5 operazioni per riga: 2 per la fattorizzazione, 3 per la sostituzione).
\end{proposizione}

\begin{keybox}[Confronto dei costi]
\begin{center}
\begin{tabular}{lcc}
\toprule
\textbf{Metodo} & \textbf{Struttura} & \textbf{Costo} \\
\midrule
LU con pivoting (SW) & Densa $N\times N$ & $O(N^3)$ \\
Thomas (spline) & Tridiagonale $(N{-}1)\times(N{-}1)$ & $O(N)$ \\
\bottomrule
\end{tabular}
\end{center}
Per $N=12$: LU $\approx 1\,152$ op., Thomas $\approx 55$ op.
Entrambi trascurabili su un calcolatore moderno.
\end{keybox}

% ============================================================
\section{Il metodo Smith--Wilson}
\label{sec:sw}
% ============================================================

\subsection{Idea e struttura}

Il metodo Smith--Wilson~\cite{smithwilson2001} costruisce la funzione di sconto
$P(t)$ su $[0,\infty)$ come combinazione di un \textbf{termine asintotico} ---
che garantisce la convergenza all'UFR --- e di $N$ \textbf{funzioni kernel},
i cui pesi $\zeta_j$ sono determinati imponendo il passaggio esatto per i dati di mercato.

\begin{keybox}[Funzione di sconto Smith--Wilson --- sez.~9.10 di \cite{eiopa2025}]
\begin{equation}
  \label{eq:sw_main}
  P(t) = e^{-\UFR_c \cdot t} + \sum_{j=1}^{N}\zeta_j\,W(t,\,u_j),
  \qquad t \ge 0,
\end{equation}
dove:
\begin{itemize}
  \item $\UFR_c = \ln(1 + \UFR_{\rm ann})$: UFR in intensit\`a continua;
  \item $u_j \in \{T_1,\ldots,T_N\}$: scadenze dei nodi di mercato (nodi kernel);
  \item $W(t,u)$: funzione kernel di Wilson (si annulla per $t\to\infty$);
  \item $\bfzeta \in \R^N$: coefficienti da determinare risolvendo il sistema lineare.
\end{itemize}
\end{keybox}

La formula vale \emph{sia} nella zona liquida ($t \le \LLP$) \emph{sia} nella zona
estrapolata ($t > \LLP$): un'unica espressione $\mathcal{C}^\infty$ per tutto $[0, 150]$.

\subsection{Funzione kernel di Wilson --- sez.~9.7 di \cite{eiopa2025}}

\begin{definizione}[Kernel di Wilson]
\begin{equation}
  \label{eq:kernel}
  W(t,u) = e^{-\UFR_c(t+u)}\Bigl[
    \alpha\cdot\min(t,u)
    - e^{-\alpha\max(t,u)}\sinh\!\bigl(\alpha\min(t,u)\bigr)
  \Bigr].
\end{equation}
\end{definizione}

\textbf{Propriet\`a} (sez.~9.9 di \cite{eiopa2025}):
\begin{enumerate}
  \item $W(t,u)=W(u,t)$ (simmetria);
  \item $W(t,u)\to 0$ esponenzialmente per $t\to\infty$, quindi $P(t)\to e^{-\UFR_c t}$;
  \item La matrice $\Wmat=[W(u_i,u_j)]$ \`e \textbf{simmetrica definita positiva}
        (conseguenza del fatto che $W$ \`e un nucleo riproducente).
\end{enumerate}

La struttura di~\eqref{eq:sw_main} --- somma di un termine asintotico e di
funzioni kernel centrate nei nodi di mercato --- non \`e una scelta arbitraria.
Come mostreremo nella Sezione~\ref{sec:variazionale}, essa \`e la soluzione ottimale
di un problema variazionale in uno spazio di Hilbert appropriato: il Teorema del
Rappresentante (Representer Theorem) garantisce che la soluzione di minima norma
ha \emph{sempre} questa forma, riducendo il problema di dimensione infinita a un
sistema lineare $N \times N$.
Lo stesso schema architetturale compare nella teoria delle
\textbf{Radial Basis Functions (RBF)}, un metodo classico di interpolazione
multidimensionale fondato sulla medesima teoria dei nuclei riproducenti.

\begin{keybox}[Analogia con le funzioni a base radiale (RBF)]
Le \textbf{RBF} (\emph{Radial Basis Functions}) costruiscono un interpolante della
stessa forma: $f(x) = \sum_j \zeta_j\,\phi(\|x-x_j\|)$, dove $\phi$ dipende
solo dalla \emph{distanza} (es.~Gaussiana $\phi(r)=e^{-r^2/\sigma^2}$, multiquadrica
$\phi(r)=\sqrt{r^2+c^2}$). Il kernel di Wilson appartiene alla stessa famiglia
concettuale (\emph{nuclei riproducenti}, RKHS), con due differenze chiave:
\begin{center}\small
\renewcommand{\arraystretch}{1.2}
\begin{tabular}{lll}
\toprule
\textbf{Caratteristica} & \textbf{RBF classica} & \textbf{Kernel di Wilson}\\
\midrule
Dipendenza & Solo da $\|x-y\|$ (radiale) & Da $t$ e $u$ separatamente\\
Termine asintotico & Assente & $e^{-\UFR_c t}$ strutturale\\
Gram matrix & Simm.\ def.\ pos. & Simm.\ def.\ pos.\\
Garanzia conv.\ all'asintoto & No (estrapolazione libera) & S\`i (verso UFR)\\
\bottomrule
\end{tabular}
\end{center}
Il fattore $e^{-\UFR_c(t+u)}$ rompe la radialit\`a: kernel con $t$ o $u$
grandi sono automaticamente smorzati, forzando la curva verso l'UFR
\emph{senza} imporre un vincolo asintotico aggiuntivo. La Figura~\ref{fig:kernel_W}
mostra la matrice di Gram $\Wmat_{ij}=W(u_i,u_j)$: simmetrica e def.\ pos.,
analogamente alla matrice di Gram RBF $[\phi(\|x_i-x_j\|)]$.
\end{keybox}

La Figura~\ref{fig:kernel} mostra $W(t,u_j)$ per diversi nodi $u_j$: il kernel
raggiunge un picco vicino a $t \approx u_j$ e decade esponenzialmente per
$t\to\infty$. Il parametro $\alpha$ controlla la velocit\`a del decadimento
(Figure~\ref{fig:kernel_alpha} e~\ref{fig:kernel_W}).

\begin{figure}[H]
\centering
\includegraphics[width=0.92\textwidth]{grafici/curva_eiopa_rfr/fig02a_kernel_wilson.pdf}
\caption{Kernel di Wilson $W(t,u_j)$ per quattro nodi rappresentativi
$u \in \{1, 5, 10, 20\}$ anni, con $\alpha=0.128$ e $\UFR_c=\ln(1.033)$.
Ogni curva si annulla per $t\to\infty$, garantendo la convergenza della
funzione di sconto al termine asintotico $e^{-\UFR_c t}$.}
\label{fig:kernel}
\end{figure}

\begin{figure}[H]
\centering
\includegraphics[width=0.92\textwidth]{grafici/curva_eiopa_rfr/fig02b_kernel_alpha_effetto.pdf}
\caption{Effetto del parametro $\alpha$ sulla forma del kernel (nodo $u=10$ fisso).
Un $\alpha$ piccolo produce un kernel pi\`u ampio e lento a decadere: la curva
impiega pi\`u tempo a convergere verso l'UFR. Un $\alpha$ grande produce un
kernel stretto e rapido.}
\label{fig:kernel_alpha}
\end{figure}

\begin{figure}[H]
\centering
\includegraphics[width=0.72\textwidth]{grafici/curva_eiopa_rfr/fig02c_matrice_kernel_W.pdf}
\caption{Matrice di Gram del kernel di Wilson: $\Wmat_{ij} = W(u_i,u_j)$ per le
12 scadenze DLT EUR, con $\alpha=0.128$. La matrice \`e simmetrica e definita
positiva (analoga alla matrice di Gram RBF): ci\`o garantisce l'unicit\`a della
soluzione del sistema di calibrazione.}
\label{fig:kernel_W}
\end{figure}

\subsection{Interpretazione variazionale --- sez.~9.6 di \cite{eiopa2025}}
\label{sec:variazionale}

La formula~\eqref{eq:sw_main} non \`e una scelta ad hoc: \`e la \textbf{soluzione
unica} di un problema di minima norma in uno spazio di Hilbert a nucleo
riproducente (RKHS, \emph{Reproducing Kernel Hilbert Space}).

\subsubsection*{Lo spazio funzionale $\mathcal{H}$ e la norma}

Fissiamo il ``prior'' asintotico $P_\infty(t)=e^{-\UFR_c t}$, che rappresenta
il comportamento atteso della curva per $t\to\infty$.
Consideriamo lo spazio $\mathcal{H}$ delle funzioni $g:[0,\infty)\to\R$
dotato di un prodotto interno costruito in modo che il kernel di Wilson
$W(t,u)$ ne sia il \textbf{nucleo riproducente}:
\[
  \langle W(\cdot,\, u),\; g \rangle_{\mathcal{H}} = g(u)
  \qquad \text{per ogni } g \in \mathcal{H},\; u \ge 0.
\]
La norma $\|g\|_{\mathcal{H}}^2 = \langle g,g\rangle_{\mathcal{H}}$ penalizza
le perturbazioni rispetto all'asintoto: una funzione ``vicina'' a $P_\infty$
ha norma piccola, una con oscillazioni elevate ha norma grande.

\subsubsection*{Il problema di minima norma con vincoli di mercato}

I tassi par di mercato $s_j$ (corretti per il CRA) definiscono $N$ vincoli
funzionali su $P$: il funzionale $\phi_j(P)$ restituisce il valore netto
dello swap $j$ attualizzato con la funzione di sconto $P$
(cfr.~eq.~\eqref{eq:tasso_par}), e deve valere 1 (swap alla pari).
Il problema variazionale di Smith--Wilson \`e:
\begin{equation}
  \label{eq:variazionale}
  \min_{P \in \mathcal{H}}\;\|P - P_\infty\|_{\mathcal{H}}^2
  \qquad\text{soggetto a:}\qquad
  \phi_j(P) = 1, \quad j=1,\ldots,N.
\end{equation}
Si minimizza la distanza da $P_\infty$ nello spazio $\mathcal{H}$ compatibilmente
con i $N$ vincoli di mercato.

\textbf{Interpretazione geometrica.}
I vincoli $\phi_j(P)=1$ definiscono un sottospazio affine di $\mathcal{H}$:
l'insieme di tutte le funzioni di sconto compatibili con i dati.
La soluzione di~\eqref{eq:variazionale} \`e la \textbf{proiezione ortogonale}
di $P_\infty$ su questo sottospazio --- la funzione del vincolo pi\`u vicina
all'asintoto nella metrica di $\mathcal{H}$.

\subsubsection*{Teorema del Rappresentante}

Il seguente risultato --- versione del Representer Theorem per RKHS
con vincoli funzionali lineari --- garantisce che la soluzione di~\eqref{eq:variazionale}
ha sempre la forma di una combinazione finita di kernel:

\begin{teorema}[Representer Theorem per Smith--Wilson]
\label{thm:representer}
La soluzione ottimale di~\eqref{eq:variazionale} \`e:
\[
  P^*(t) = P_\infty(t) + \sum_{j=1}^{N} \zeta_j\,W(t,\,u_j),
\]
dove i coefficienti $\bfzeta = (\zeta_1,\ldots,\zeta_N)^\top$ sono univocamente
determinati dai vincoli di mercato.
\end{teorema}

Nonostante~\eqref{eq:variazionale} sia un problema in dimensione infinita,
la soluzione risiede nel sottospazio finito-dimensionale generato dai kernel
$\{W(\cdot,u_j)\}_{j=1}^N$: la complessit\`a si riduce a un sistema lineare $N\times N$.

\subsubsection*{Dal problema variazionale al sistema lineare}

Sostituendo $P^*(t)$ nel vincolo $\phi_j(P^*) = 1$ e sfruttando la
linearit\`a di $\phi_j$:
\[
  \phi_j(P_\infty) + \sum_{k=1}^{N} \zeta_k\,\phi_j\!\bigl(W(\cdot,u_k)\bigr) = 1.
\]
Definendo $H_{jk} \defeq \phi_j(W(\cdot,u_k))$ e
$\mu_{Q,j} \defeq \phi_j(P_\infty)$, si ottiene esattamente il sistema
della Sezione~\ref{sec:sistema}:
\[
  \Hmat\,\bfzeta = \mathbf{1} - \bfmu_Q.
\]
La struttura del problema variazionale \`e quindi la \emph{radice matematica}
del sistema lineare denso: ogni elemento $H_{jk}$ misura quanto il kernel
nel nodo $u_k$ contribuisce al valore dello swap $j$.

\begin{keybox}[Parallelo con la spline cubica naturale]
La spline cubica naturale risolve:
\[
  \min_{S \in C^2[a,b]}\;\int_a^b \bigl[S''(t)\bigr]^2\,\mathrm{d}t
  \quad\text{s.t.}\quad S(t_i) = y_i,\quad i=1,\ldots,n.
\]
Smith--Wilson \`e il diretto analogo con due varianti essenziali:
\begin{enumerate}
  \item La \textbf{norma} \`e quella di $\mathcal{H}$, definita dal kernel
        $W(t,u)$ invece di $\int[S'']^2$: questo incorpora strutturalmente
        la convergenza all'UFR, senza vincolo esplicito aggiuntivo.
  \item I \textbf{vincoli} sono funzionali swap $\phi_j(P)=1$ invece di
        valutazioni puntali $S(t_i)=y_i$: riflettono la struttura degli
        strumenti osservabili sul mercato.
\end{enumerate}
In entrambi i casi il Representer Theorem garantisce che la soluzione ottimale
\`e una combinazione lineare di kernel centrati nei nodi dei vincoli,
i cui pesi si trovano risolvendo un sistema lineare $N\times N$.
\end{keybox}

\subsection{Parametri chiave per EUR 2025}

\begin{center}
\renewcommand{\arraystretch}{1.25}
\begin{tabular}{lll}
\toprule
\textbf{Parametro} & \textbf{Valore EUR 2025} & \textbf{Riferimento} \\
\midrule
$\UFR_{\rm ann}$ & $3.30\%$ (annuo composto) & sez.~9.3, Annex~E \\
$\UFR_c$ & $\ln(1.033) \approx 3.2466\%$ (continuo) & conversione \\
$\LLP$ & 20 anni & sez.~3.7 (residual volume) \\
$\CP$ & $\max(\LLP+40,\,60) = 60$ anni & sez.~9.4 \\
$\CRA$ & 10 bps & sez.~7.3, Situazione~1 \\
\bottomrule
\end{tabular}
\end{center}

\begin{warningbox}
\textbf{Attenzione alla convenzione dell'UFR.}
$\UFR_{\rm ann} = 3.30\%$ \`e in convenzione annua composta ($1+r$), \emph{non}
in capitalizzazione continua ($e^r$). Nelle formule Smith--Wilson si usa
sempre l'intensit\`a continua:
\[
  \UFR_c = \ln(1 + 0.033) \approx 0.032466 \ne 0.0330.
\]
L'errore di usare $0.0330$ al posto di $0.032466$ produce uno scostamento
di $\sim 5$ bps sulla curva a lungo termine --- significativo nel contesto
regolamentare.
\end{warningbox}

% ============================================================
\section{Il sistema lineare: calibrazione su swap par}
\label{sec:sistema}
% ============================================================

\subsection{Dati di input: OIS e CRA}

Per l'euro, EIOPA utilizza come input i tassi par degli OIS su \texteuro{}STR
(sez.~5.3, Annex~B di~\cite{eiopa2025}). Un OIS par a scadenza $T_k$ con cedola
annua e nominale 1 vale esattamente 1 sul mercato:
\begin{equation}
  \label{eq:ois_par}
  s_k \sum_{j=1}^{T_k} P(0,j) + P(0,T_k) = 1,
\end{equation}
dove $P(0,t)$ \`e il fattore di sconto (incognita).

Il \textbf{Credit Risk Adjustment} (sez.~7.3) viene sottratto prima della
calibrazione: $r_k = s_k - \CRA$, con $\CRA = 10$ bps per EUR 2025.

\begin{table}[H]
\centering\small
\caption{Tassi OIS \texteuro{}STR par (illustrativi, maggio 2025) e tassi after-CRA.}
\label{tab:ois_eur}
\begin{tabular}{ccccc}
\toprule
$T_k$ (anni) & $s_k$ (\%) & $r_k = s_k - 0.10$ (\%) & DLT & \\
\midrule
\rowcolor{rowtint}1  & 3.742 & 3.642 & \checkmark & \\
2  & 2.851 & 2.751 & \checkmark & \\
\rowcolor{rowtint}3  & 2.577 & 2.477 & \checkmark & \\
4  & 2.468 & 2.368 & \checkmark & \\
\rowcolor{rowtint}5  & 2.413 & 2.313 & \checkmark & \\
6  & 2.398 & 2.298 & \checkmark & \\
\rowcolor{rowtint}7  & 2.393 & 2.293 & \checkmark & \\
8  & 2.400 & 2.300 & \checkmark & \\
\rowcolor{rowtint}10 & 2.470 & 2.370 & \checkmark & \\
12 & 2.510 & 2.410 & \checkmark & \\
\rowcolor{rowtint}15 & 2.555 & 2.455 & \checkmark & \\
20 & 2.581 & 2.481 & \checkmark & \textbf{LLP}\\
\bottomrule
\end{tabular}
\end{table}

\begin{figure}[H]
\centering
\includegraphics[width=0.92\textwidth]{grafici/curva_eiopa_rfr/fig01_input_ois_par.pdf}
\caption{Tassi OIS \texteuro{}STR par (lordi) e after-CRA per le 12 scadenze DLT EUR.
Il tasso a 1 anno \`e significativamente pi\`u alto: riflette la politica monetaria
restrittiva corrente. La curva si appiattisce oltre i 5 anni.}
\label{fig:input}
\end{figure}

\subsection{Derivazione del sistema --- sez.~9.15 di \cite{eiopa2025}}

EIOPA applica Smith--Wilson \emph{direttamente} ai tassi par degli swap, senza
alcun passo intermedio. Ogni swap par a scadenza $T_i$ con tasso $r_i$ ha
valore di mercato 1 (condizione~\eqref{eq:ois_par}), con flussi di cassa:
\[
  c_{ij} = \begin{cases}
    r_i & j = 1, \ldots, T_i - 1, \\
    1 + r_i & j = T_i.
  \end{cases}
\]

Sostituendo~\eqref{eq:sw_main} in $\sum_j c_{ij} P(j) = 1$ e separando
i termini noti e incogniti:
\[
  \underbrace{\sum_{j=1}^{T_i} c_{ij}\,e^{-\UFR_c\,j}}_{\mu_{Q,i}}
  + \sum_{k=1}^{N}\zeta_k \underbrace{\sum_{j=1}^{T_i} c_{ij}\,W(j,\,u_k)}_{H_{ik}}
  = 1.
\]

\begin{keybox}[Sistema lineare di calibrazione EIOPA --- sez.~9.15]
\begin{equation}
  \label{eq:H_system}
  \Hmat\,\bfzeta = \mathbf{1} - \bfmu_Q,
\end{equation}
con:
\begin{align}
  \mu_{Q,i} &= \sum_{j=1}^{T_i} c_{ij}\,e^{-\UFR_c\,j}
               \quad\text{(contributo del termine asintotico)},
               \label{eq:mu_Q}\\
  H_{ik}    &= \sum_{j=1}^{T_i} c_{ij}\,W(j,\,u_k)
               \quad\text{(matrice $N\times N$ densa)}.
               \label{eq:H_ik}
\end{align}
\end{keybox}

\begin{esempio}[Riga $H_{ik}$ per $T_i = 3$, $r_i = 2.477\%$]
I flussi sono $c_{i1}=0.02477$, $c_{i2}=0.02477$, $c_{i3}=1.02477$.
\[
  H_{ik} = 0.02477\,W(1,u_k) + 0.02477\,W(2,u_k) + 1.02477\,W(3,u_k).
\]
Ogni riga \`e la somma pesata (con pesi dati dai flussi cedola) dei valori
del kernel. La matrice $\Hmat$ \`e \textbf{densa} e in generale \textbf{non simmetrica}
(a differenza della matrice $\Wmat$ pura).
\end{esempio}

\subsection{Risoluzione: LU con pivoting parziale}

Il sistema~\eqref{eq:H_system} si risolve con la fattorizzazione LU presentata
nella Sezione~\ref{sec:lu}:
\[
  P_\pi \Hmat = LU, \qquad
  Ly = P_\pi(\mathbf{1}-\bfmu_Q), \qquad
  U\bfzeta = y.
\]
Costo: $\frac{2}{3}N^3 + O(N^2) \approx 1\,152$ operazioni per $N=12$.

\begin{figure}[H]
\centering
\includegraphics[width=0.65\textwidth]{grafici/curva_eiopa_rfr/fig03a_matrice_H.pdf}
\caption{Struttura della matrice $\Hmat$ (heatmap): densa, non simmetrica.
Le righe corrispondono alle scadenze $T_i$ degli swap, le colonne ai nodi kernel $u_k$.}
\label{fig:matrice_H}
\end{figure}

\begin{figure}[H]
\centering
\includegraphics[width=0.85\textwidth]{grafici/curva_eiopa_rfr/fig03b_coefficienti_zeta.pdf}
\caption{Coefficienti $\zeta_k$ ottenuti dalla risoluzione del sistema
$\Hmat\bfzeta = \mathbf{b}$ con $\alpha=0.128$. I $\zeta_k$ traducono i vincoli
di mercato in ``correzioni'' alla curva asintotica $e^{-\UFR_c t}$.}
\label{fig:zeta}
\end{figure}

\subsection{Condizionamento della matrice $\Hmat$}

\begin{figure}[H]
\centering
\includegraphics[width=0.85\textwidth]{grafici/curva_eiopa_rfr/fig03c_condizionamento_vs_alpha.pdf}
\caption{Numero di condizionamento $\kappa_2(\Hmat)$ in funzione di $\alpha$.
Per $\alpha$ piccolo le colonne di $\Hmat$ diventano quasi proporzionali
(il kernel $W(j,u_k)$ dipende poco da $u_k$), aumentando $\kappa_2$.
Per $\alpha \in [0.10, 0.50]$ il sistema \`e ben condizionato.}
\label{fig:cond}
\end{figure}

\subsection{Valutazione della curva}

Una volta risolto~\eqref{eq:H_system} e ottenuto $\bfzeta$, la formula~\eqref{eq:sw_main}
fornisce $P(t)$ per ogni $t \ge 0$:
\begin{align}
  r(0,t) &= -\frac{\ln P(t)}{t}
             \quad\text{(tasso spot continuo)},
             \label{eq:spot_sw}\\[4pt]
  f(0,t) &= -\frac{\mathrm{d}}{\mathrm{d}t}\ln P(t)
           = \UFR_c - \frac{\displaystyle\sum_{j=1}^{N}\zeta_j\,\partial_t W(t,u_j)}{P(t)}
             \quad\text{(tasso forward istantaneo)}.
             \label{eq:fwd_sw}
\end{align}

Per $t\to\infty$ il numeratore di~\eqref{eq:fwd_sw} va a zero (i kernel
si annullano esponenzialmente), quindi $f(0,t)\to\UFR_c$: la convergenza
\`e \textbf{strutturale}, incorporata nella formula.

% ============================================================
\section{Determinazione del parametro $\alpha$}
\label{sec:alpha}
% ============================================================

\subsection{L'Ultimate Forward Rate}

L'UFR \`e il tasso forward di lungo periodo stimato come somma di
tasso reale atteso e inflazione attesa (sez.~9.3, Annex~E di~\cite{eiopa2025}):
\[
  \UFR_{\rm ann} = r_{\rm reale}^{\rm att} + \pi^{\rm att}
  \approx 1.30\% + 2.00\% = 3.30\%.
\]

\subsection{Criterio EIOPA per $\alpha$ --- sez.~9.4}

La velocit\`a di convergenza verso l'UFR \`e controllata da $\alpha$.
Il criterio: $\alpha^*$ \`e il \emph{minimo} valore in $[0.05, 1.0]$
tale che il forward al Convergence Point si discosti dall'UFR di meno di 1 bp:
\begin{equation}
  \label{eq:alpha_crit}
  g(\alpha) \;\defeq\; f(0,\CP;\,\alpha) - \UFR_c, \qquad
  \abs{g(\alpha^*)} \le 10^{-4}.
\end{equation}

La funzione $g(\alpha)$ dipende da $\alpha$ in modo non lineare: per ogni valore
di $\alpha$ occorre risolvere il sistema $\Hmat(\alpha)\bfzeta = \mathbf{b}$ e
poi valutare il forward al CP.

\begin{proposizione}[Monotonia e applicabilit\`a della bisezione]
$g(\alpha)$ \`e strettamente monotona decrescente (in valore assoluto) su
$[0.05, 1.0]$: al crescere di $\alpha$, la convergenza verso l'UFR \`e pi\`u
rapida e il forward al CP \`e pi\`u vicino all'UFR. Il metodo di bisezione
della Sezione~\ref{sec:bisezione} \`e applicabile con convergenza garantita.
\end{proposizione}

\subsection{Algoritmo: bisezione per $\alpha^*$}

\begin{algorithm}[H]
\caption{Bisezione per $\alpha^*$ (alternativa alla scanning procedure EIOPA)}
\label{alg:alpha}
\DontPrintSemicolon
\SetKwInOut{Input}{Input}\SetKwInOut{Output}{Output}
\Input{$\{T_k\}$, $\{r_k\}$, $\UFR_c$, $\CP$, $\varepsilon_\alpha=10^{-6}$, $n_{\max}=60$}
\Output{$\alpha^*$ = valore minimo che soddisfa~\eqref{eq:alpha_crit}}
$a\leftarrow 0.05$;\quad $b\leftarrow 1.0$\;
\For{$k=1,\ldots,n_{\max}$}{
  $\alpha_{\rm mid}\leftarrow (a+b)/2$\;
  Risolvi $\Hmat(\alpha_{\rm mid})\bfzeta = \mathbf{b}$ $\longrightarrow$ $\bfzeta$\;
  Calcola $f_{\rm mid} = f(0,\CP;\alpha_{\rm mid})$ da~\eqref{eq:fwd_sw}\;
  \eIf{$\abs{f_{\rm mid}-\UFR_c}\le 10^{-4}$}{
    $b\leftarrow\alpha_{\rm mid}$ \tcp*{soddisfatto: cercare $\alpha$ pi\`u piccolo}
  }{
    $a\leftarrow\alpha_{\rm mid}$ \tcp*{non soddisfatto: $\alpha$ troppo piccolo}
  }
  \If{$b-a<\varepsilon_\alpha$}{\textbf{break}\;}
}
\Return{$b$}\;
\end{algorithm}

\paragraph{Analisi dell'errore.} Applicando~\eqref{eq:bisez_err} all'intervallo
$[0.05,\,1.0]$ (ampiezza $\ell = 0.95$):
\[
  \abs{\alpha^* - \alpha_n} \le \frac{0.95}{2^n}.
\]
Per $\varepsilon_\alpha = 10^{-6}$: $n \ge \lceil\log_2(0.95/10^{-6})\rceil = 20$
iterazioni. Per EUR 2025: $\alpha^* \approx 0.128$.

\paragraph{Confronto con la scanning procedure EIOPA.}
EIOPA usa una griglia densa con passo $10^{-6}$ su $[0.05, 1.0]$:
circa $950\,000$ valutazioni di $g$. La bisezione ne richiede $\approx 20$.
La bisezione \`e $\sim 47\,000$ volte pi\`u efficiente.

\paragraph{Costo totale.} Ogni valutazione di $g$ richiede la risoluzione di
$\Hmat\bfzeta = \mathbf{b}$: costo $O(N^3)$. Totale bisezione:
$20 \times 12^3 \approx 34\,560$ operazioni.

\begin{figure}[H]
\centering
\includegraphics[width=0.92\textwidth]{grafici/curva_eiopa_rfr/fig04b_funzione_g_alpha.pdf}
\caption{La funzione $g(\alpha) = f(0,\CP;\alpha) - \UFR_c$ in bps.
La zona verde indica la regione $|g|\le 1$ bp dove il criterio EIOPA \`e soddisfatto.
$\alpha^*$ \`e il minimo valore in questa zona.}
\label{fig:g_alpha}
\end{figure}

\begin{figure}[H]
\centering
\includegraphics[width=0.92\textwidth]{grafici/curva_eiopa_rfr/fig04a_bisezione_convergenza.pdf}
\caption{Convergenza della bisezione: ampiezza dell'intervallo $[a_n, b_n]$
in funzione dell'iterazione (scala logaritmica). La convergenza \`e
esponenziale: ogni iterazione dimezza l'intervallo.}
\label{fig:bisez_conv}
\end{figure}

\begin{keybox}[Bilancio computazionale: pipeline completa]
\begin{center}\small
\renewcommand{\arraystretch}{1.25}
\begin{tabular}{lrrl}
\toprule
\textbf{Fase} & \textbf{LU calls} & \textbf{Op.\ floating-point} & \textbf{Nota}\\
\midrule
Bisezione ($n=20$ iter.) & 20 & $20\times 1\,152 \approx 23\,000$ & $\frac{2}{3}N^3$, $N=12$\\
Valutazione $P(t)$, $f(t)$ & 0 & $O(N\cdot M)\approx 2\,400$ & $M=200$ punti\\
\midrule
\textbf{Totale pipeline} & 20 & $\sim 25\,000$ & $\sim 25\,\mu$s su CPU moderna\\
\midrule
Scanning EIOPA (passo $10^{-6}$) & $950\,000$ & $\sim 10^9$ & $\approx\!50\,000\times$ pi\`u lento\\
\bottomrule
\end{tabular}
\end{center}
La differenza \`e irrilevante per un singolo calcolo mensile, ma diventa
sostanziale quando la curva viene ricalcolata su migliaia di scenari Monte
Carlo (stress test SCR, ORSA): la bisezione permette di mantenere il calcolo
nel dominio del tempo reale.
\end{keybox}

% ============================================================
\section{Risultati: la curva EIOPA completa}
\label{sec:risultati}
% ============================================================

\subsection{Curva spot e forward}

Eseguendo la pipeline completa con $\alpha^*$ trovato dalla bisezione si
ottiene la curva $P(t)$ per $t \in [0, 150]$, da cui spot e forward.

\begin{figure}[H]
\centering
\includegraphics[width=0.95\textwidth]{grafici/curva_eiopa_rfr/fig05a_curva_spot_forward.pdf}
\caption{Curva EIOPA EUR completa: tasso spot $r(0,T)$ (blu continuo) e forward
istantaneo $f(0,T)$ (rosso tratteggiato). I punti neri sono gli spot ai nodi
di mercato. Le linee verticali indicano LLP~(20a) e CP~(60a). Oltre il CP,
il forward converge all'UFR~$=3.30\%$.}
\label{fig:curva_completa}
\end{figure}

\begin{figure}[H]
\centering
\includegraphics[width=0.85\textwidth]{grafici/curva_eiopa_rfr/fig05b_convergenza_ufr.pdf}
\caption{Distanza $|f(0,T) - \UFR_c|$ in bps (scala logaritmica).
Al Convergence Point ($T=60$) il forward \`e per costruzione entro 1~bp dall'UFR.
La convergenza \`e esponenziale, governata dal termine $e^{-\alpha T}$ del kernel.}
\label{fig:conv_ufr}
\end{figure}

\begin{figure}[H]
\centering
\includegraphics[width=0.85\textwidth]{grafici/curva_eiopa_rfr/fig05c_fattori_sconto.pdf}
\caption{Fattori di sconto $P(0,T)$ per $T=1,\ldots,100$ anni.
La curva puntinata viola \`e il limite asintotico $e^{-\UFR_c T}$: per
$T$ grande, $P(T)$ converge a questa curva (i kernel si annullano).}
\label{fig:discount}
\end{figure}

\subsection{Verifica: replica dei par swap}

La calibrazione \`e corretta se i fattori di sconto ottenuti replicano
esattamente i par swap di input. Per ogni strumento $i$:
\[
  r_i \sum_{j=1}^{T_i} P(0,j) + P(0,T_i) \stackrel{?}{=} 1.
\]
Con il codice R si verifica che l'errore massimo \`e dell'ordine di $10^{-15}$
(precisione macchina), confermando la correttezza della risoluzione del sistema.

% ============================================================
\section{Confronto: spline cubiche vs Smith--Wilson}
\label{sec:confronto}
% ============================================================

Le spline cubiche non sono nella pipeline EIOPA (sez.~9.1 di~\cite{eiopa2025}),
ma il loro confronto con Smith--Wilson \`e illuminante per il calcolo numerico:
i due metodi risolvono lo stesso problema di interpolazione con strutture
algebriche molto diverse.

\subsection{Sistema tridiagonale per le spline}

Data la funzione $\phi(T) = -\ln P(0,T)$ ai nodi $T_0 < \cdots < T_N$
con passi $h_j = T_{j+1}-T_j$, le derivate seconde $M_j = S''(T_j)$ soddisfano
il sistema tridiagonale:
\begin{equation}
  \label{eq:spline_sys}
  h_{j-1}M_{j-1} + 2(h_{j-1}+h_j)M_j + h_j M_{j+1}
  = 6\!\left(\frac{\phi_{j+1}-\phi_j}{h_j} - \frac{\phi_j-\phi_{j-1}}{h_{j-1}}\right).
\end{equation}

In forma matriciale $\mathbf{A}\mathbf{M} = \mathbf{d}$ con $\mathbf{A}$
tridiagonale, simmetrica, a dominanza diagonale stretta. Si risolve con
l'algoritmo di Thomas (Proposizione~\ref{prop:thomas}) in $O(N)$.

\subsection{Confronto strutturale e limiti della spline}

\paragraph{Perch\`e il sistema SW \`e denso e quello delle spline \`e tridiagonale?}
La risposta sta nella \emph{localit\`a} delle funzioni base:
\begin{itemize}
  \item Le spline cubiche usano polinomi a \textbf{supporto locale}: il tratto
        su $[T_j, T_{j+1}]$ dipende solo dai nodi adiacenti. Ogni condizione di
        raccordo coinvolge al pi\`u 3 nodi consecutivi $\Rightarrow$ matrice
        \textbf{tridiagonale}, algoritmo di Thomas $O(N)$.
  \item Il kernel di Wilson $W(t,u_k)$ ha \textbf{supporto globale}: ogni nodo
        $u_k$ influenza la curva su tutto $[0,\infty)$. Ogni condizione di
        mercato coinvolge \emph{tutti} i nodi kernel contemporaneamente
        $\Rightarrow$ matrice $\Hmat$ \textbf{densa}, fattorizzazione LU $O(N^3)$.
\end{itemize}
Il prezzo della garanzia strutturale di convergenza all'UFR (kernel globale)
\`e una matrice densa. Le spline, non avendo questo requisito, sfruttano la
struttura locale e ottengono un sistema tridiagonale molto pi\`u economico.

\begin{figure}[H]
\centering
\includegraphics[width=0.92\textwidth]{grafici/curva_eiopa_rfr/fig06_confronto_spline_sw.pdf}
\caption{Confronto interpolazione/estrapolazione: spline cubica naturale vs
Smith--Wilson. Nella zona liquida ($T\le 20$) entrambe passano per i nodi.
Oltre il LLP la spline estrapolizza linearmente (senza convergenza a UFR),
mentre SW converge strutturalmente.}
\label{fig:confronto}
\end{figure}

\begin{center}\small
\renewcommand{\arraystretch}{1.25}
\begin{tabular}{lcc}
\toprule
\textbf{Propriet\`a} & \textbf{Spline cubica} & \textbf{Smith--Wilson} \\
\midrule
Struttura del sistema & Tridiagonale $\mathbf{A}$ ($N-1$ variabili) & Densa $\Hmat$ ($N$ variabili) \\
Algoritmo di soluzione & Thomas: $O(N)$ & LU: $O(N^3)$ \\
Interpolazione & \checkmark & \checkmark \\
Estrapolazione con conv.\ a UFR & $\times$ & \checkmark \\
Formula unica $[0,\infty)$ & $\times$ & \checkmark \\
Regolarit\`a globale & $\mathcal{C}^2$ & $\mathcal{C}^\infty$ \\
\bottomrule
\end{tabular}
\end{center}

% ============================================================
\section{Analisi di sensitivit\`a}
\label{sec:sensitivita}
% ============================================================

\subsection{Sensitivit\`a all'UFR}

\begin{figure}[H]
\centering
\includegraphics[width=0.92\textwidth]{grafici/curva_eiopa_rfr/fig07a_sensitivita_ufr.pdf}
\caption{Curva spot per tre valori dell'UFR: 3.15\%, 3.30\%, 3.45\%.
Nella zona liquida ($T\le 20$) le curve sono indistinguibili;
oltre il LLP divergono verso il rispettivo UFR.}
\label{fig:sens_ufr}
\end{figure}

\subsection{Sensitivit\`a al parametro $\alpha$}

\begin{figure}[H]
\centering
\includegraphics[width=0.92\textwidth]{grafici/curva_eiopa_rfr/fig07b_sensitivita_alpha.pdf}
\caption{Forward istantaneo per tre valori di $\alpha$ (UFR fisso a 3.30\%).
Un $\alpha$ pi\`u grande produce una convergenza pi\`u rapida verso l'UFR,
con la curva forward che ``raggiunge'' l'asintoto prima.}
\label{fig:sens_alpha}
\end{figure}

% ============================================================
\section{Implementazione in R}
\label{sec:codice}
% ============================================================

Il codice completo \`e nel file \texttt{R/curva\_eiopa\_rfr\_CN.R}.
Qui presentiamo le funzioni principali con commenti dettagliati.

\subsection{Setup e dati}

\begin{codebox}[Parametri EIOPA e dati di input]
\begin{Verbatim}[fontsize=\small]
# Parametri EIOPA EUR 2025 (sez. 9.3, 9.4)
UFR_ann <- 0.0330                   # 3.30% annuo composto
UFR_c   <- log(1 + UFR_ann)        # intensità continua ~ 3.2466%
LLP     <- 20                       # Last Liquid Point EUR
CP      <- max(LLP + 40, 60)       # Convergence Point = 60 anni
CRA     <- 0.0010                   # 10 bps

# Scadenze DLT e tassi par OIS €STR (illustrativi, maggio 2025)
T_mkt <- c(1, 2, 3, 4, 5, 6, 7, 8, 10, 12, 15, 20)
s_mkt <- c(3.742, 2.851, 2.577, 2.468, 2.413, 2.398,
           2.393, 2.400, 2.470, 2.510, 2.555, 2.581) / 100
r_mkt <- s_mkt - CRA   # tassi after-CRA
\end{Verbatim}
\end{codebox}

\subsection{Kernel di Wilson}

\begin{codebox}[Implementazione del kernel --- eq.~\eqref{eq:kernel}]
\begin{Verbatim}[fontsize=\small]
W_kernel <- function(t, u, ufr, alpha) {
  tmin <- pmin(t, u)
  tmax <- pmax(t, u)
  exp(-ufr * (t + u)) * (alpha * tmin
    - exp(-alpha * tmax) * sinh(alpha * tmin))
}
\end{Verbatim}
\end{codebox}

\subsection{Costruzione della matrice $\Hmat$ e calibrazione}

\begin{codebox}[Calibrazione SW su swap par --- sez.~9.15 EIOPA]
\begin{Verbatim}[fontsize=\small]
sw_calibrate <- function(maturities, par_rates, ufr, alpha) {
  N <- length(maturities)
  H    <- matrix(0, N, N)
  mu_Q <- numeric(N)
  
  for (i in seq_len(N)) {
    Ti <- maturities[i]; ri <- par_rates[i]
    for (j in seq_len(Ti)) {
      c_ij <- if (j < Ti) ri else (1 + ri)
      mu_Q[i] <- mu_Q[i] + c_ij * exp(-ufr * j)
      for (k in seq_len(N))
        H[i, k] <- H[i, k] + c_ij * W_kernel(j, maturities[k], ufr, alpha)
    }
  }
  # Sistema lineare denso: LU con pivoting (LAPACK dgesv via R)
  zeta <- solve(H, 1 - mu_Q)
  list(zeta = zeta, H = H, b = 1 - mu_Q, mu_Q = mu_Q)
}
\end{Verbatim}
\end{codebox}

\begin{osservazione}
La funzione \texttt{solve()} di R invoca la routine LAPACK \texttt{dgesv}:
fattorizzazione LU con pivoting parziale, esattamente il metodo della
Sezione~\ref{sec:lu}. Il costo per $N=12$ \`e trascurabile ($\sim 1\,\mu$s).
\end{osservazione}

\subsection{Bisezione per $\alpha^*$}

\begin{codebox}[Bisezione --- Algoritmo~\ref{alg:alpha} in R]
\begin{Verbatim}[fontsize=\small]
find_alpha <- function(maturities, par_rates, ufr, llp,
                       tol_fwd = 1e-4, tol_alpha = 1e-6, n_max = 60) {
  cp <- max(llp + 40, 60)
  a <- 0.05; b <- 1.0
  
  for (k in seq_len(n_max)) {
    am <- 0.5 * (a + b)
    cal <- sw_calibrate(maturities, par_rates, ufr, am)
    fwd_cp <- sw_forward(cp, maturities, cal$zeta, ufr, am)
    
    if (abs(fwd_cp - ufr) <= tol_fwd) {
      b <- am   # soddisfatto: cercare alpha più piccolo
    } else {
      a <- am   # non soddisfatto
    }
    if ((b - a) < tol_alpha) break
  }
  b
}
\end{Verbatim}
\end{codebox}

\subsection{Valutazione della curva}

\begin{codebox}[Fattore di sconto e forward]
\begin{Verbatim}[fontsize=\small]
sw_discount <- function(t, maturities, zeta, ufr, alpha) {
  P <- exp(-ufr * t)
  for (k in seq_along(maturities))
    P <- P + zeta[k] * W_kernel(t, maturities[k], ufr, alpha)
  P
}

sw_forward <- function(t, maturities, zeta, ufr, alpha, h = 1e-7) {
  Pp <- sw_discount(t + h, maturities, zeta, ufr, alpha)
  Pm <- sw_discount(t - h, maturities, zeta, ufr, alpha)
  -(log(Pp) - log(Pm)) / (2 * h)
}
\end{Verbatim}
\end{codebox}

\begin{osservazione}
Il forward \`e calcolato con una \textbf{derivata centrata}
($h = 10^{-7}$). L'errore di troncamento \`e $O(h^2) \approx 10^{-14}$,
ben sotto la precisione richiesta (1~bp $= 10^{-4}$). In alternativa,
si pu\`o usare la derivata analitica $\partial_t W(t,u)$ fornita nella sez.~9.12
di~\cite{eiopa2025}.
\end{osservazione}

% ============================================================
\section*{Riepilogo degli algoritmi}
\addcontentsline{toc}{section}{Riepilogo degli algoritmi}
% ============================================================

\begin{center}\small
\renewcommand{\arraystretch}{1.35}
\begin{tabular}{llllc}
\toprule
\textbf{Fase} &
\textbf{Problema numerico} &
\textbf{Algoritmo} &
\textbf{Costo} &
\textbf{Rif.\ EIOPA}\\
\midrule
Calibrazione SW &
$\Hmat\bfzeta = \mathbf{b}$ (densa, $N\times N$) &
LU con pivoting &
$O(2N^3/3)$ &
sez.~9.15\\
Parametro $\alpha$ &
$g(\alpha^*)=0$ (non lineare) &
Bisezione &
$O(n\cdot N^3)$ &
sez.~9.4\\
\midrule
\emph{Confronto:} Spline &
$\mathbf{A}\mathbf{M}=\mathbf{d}$ (tridiag.) &
Thomas &
$O(N)$ &
---\\
\bottomrule
\end{tabular}
\end{center}

% ============================================================
\section{Esercizi}
\label{sec:esercizi}
% ============================================================

\begin{esercizio}[Costruzione della matrice $\Hmat$ a mano]
Considerare $N=3$ swap par a scadenze $T=(1,2,3)$ anni con tassi
$r=(3.50,\, 3.00,\, 2.75)\%$ after-CRA, $\UFR_c=\ln(1.033)$, $\alpha=0.15$.
\begin{enumerate}[label=(\alph*)]
  \item Scrivete esplicitamente i flussi $c_{ij}$ per ciascuno swap.
  \item Calcolate $\mu_{Q,i}$ per $i=1,2,3$.
  \item Calcolate $W(j, u_k)$ per $j,k \in \{1,2,3\}$ con la formula~\eqref{eq:kernel}.
  \item Costruite $\Hmat_{3\times 3}$ e il vettore $\mathbf{1}-\bfmu_Q$.
  \item Risolvete il sistema con eliminazione gaussiana e trovate $\bfzeta$.
  \item Verificate: $P(0,1),P(0,2),P(0,3)$ soddisfano la condizione~\eqref{eq:ois_par}?
\end{enumerate}
\end{esercizio}

\begin{esercizio}[Numero di condizionamento e stabilit\`a]
Usare i dati della Tabella~\ref{tab:ois_eur} con $\alpha=0.128$.
\begin{enumerate}[label=(\alph*)]
  \item Calcolate $\kappa_2(\Hmat)$ con R (\texttt{kappa(H, exact=TRUE)}).
  \item Perturbate il tasso $r_{20}$ di $\pm 5$ bps. Di quanto cambia $\bfzeta$?
        La stima $\|\Delta\bfzeta\| \le \kappa_2(\Hmat)\|\Delta\mathbf{b}\|/\|\mathbf{b}\|$
        \`e soddisfatta?
  \item Ripetete con $\alpha=0.05$ e $\alpha=0.50$. Come cambia $\kappa_2$?
  \item Tracciate $\kappa_2(\Hmat(\alpha))$ per $\alpha \in [0.05, 0.50]$ e confrontate
        con la Figura~\ref{fig:cond}.
\end{enumerate}
\end{esercizio}

\begin{esercizio}[Bisezione per $\alpha$]
$\LLP=20$, $\UFR_c=\ln(1.033)$, $\CP=60$.
\begin{enumerate}[label=(\alph*)]
  \item Implementate $g(\alpha) = f(0,60;\alpha) - \UFR_c$ e tracciate il grafico
        per $\alpha \in [0.05, 0.50]$. Confrontate con la Figura~\ref{fig:g_alpha}.
  \item Applicate la bisezione su $[0.05, 1.0]$ con $\varepsilon=10^{-6}$.
        Quante iterazioni servono? Confrontate con la stima~\eqref{eq:bisez_iter}.
  \item Implementate Newton--Raphson con derivata numerica centrata
        ($h=10^{-4}$). Quante iterazioni? Confrontate con la bisezione.
  \item (Facoltativo) Implementate la scanning procedure EIOPA (griglia con passo $10^{-6}$).
        Quante valutazioni di $g$ richiede? Confrontate i tempi con \texttt{system.time()}.
\end{enumerate}
\end{esercizio}

\begin{esercizio}[Convergenza verso l'UFR]
Con $\alpha^*$ trovato nell'esercizio precedente:
\begin{enumerate}[label=(\alph*)]
  \item Calcolate $f(0,T)$ per $T=20,30,40,50,60,70,80,100,150$.
  \item In quale anno $T^*$ si ha $|f(0,T^*) - \UFR_c| < 0.1$ bps per la prima volta?
  \item Ripetete con $\alpha = \alpha^* - 0.05$ e $\alpha = \alpha^* + 0.05$.
        Come cambia il profilo di convergenza?
  \item Spiegate perch\'e ridurre $\alpha$ richiede $T$ pi\`u grandi per
        raggiungere $\UFR_c$, usando la formula del kernel~\eqref{eq:kernel}.
\end{enumerate}
\end{esercizio}

\begin{esercizio}[Progetto: replica della curva EIOPA EUR]
\begin{enumerate}[label=(\alph*)]
  \item Eseguite la pipeline completa con i dati della Tabella~\ref{tab:ois_eur}
        e producete la curva spot 1--150 anni.
  \item Producete un grafico con: spot, forward istantaneo, nodi di mercato, LLP, CP, UFR.
  \item Sensitivit\`a: variate $\UFR_{\rm ann} \in \{3.15\%,\,3.30\%,\,3.45\%\}$ e
        confrontate le curve.
  \item Sensitivit\`a: variate $\alpha \in \{\alpha^*-0.02,\,\alpha^*,\,\alpha^*+0.02\}$
        e confrontate.
  \item \emph{Sfida}: scaricate il file Excel mensile EIOPA da
        \url{https://www.eiopa.europa.eu/tools-and-data} e confrontate
        la vostra curva spot con quella pubblicata. L'errore massimo deve essere
        $<0.1$ bps.
\end{enumerate}
\end{esercizio}

% ============================================================
\renewcommand{\refname}{Riferimenti bibliografici}
\begin{thebibliography}{9}

\bibitem{eiopa2025}
EIOPA (2025).
\textit{Technical Documentation of the Methodology to Derive EIOPA's
Risk-Free Interest Rate Term Structures}.
EIOPA-BoS-25-599, Dicembre 2025.
\url{https://www.eiopa.europa.eu}

\bibitem{eiopa2025b}
EIOPA (2025).
\textit{Updated Technical RFR Documentation Applicable as of 30 June 2025}.
EIOPA-BoS-25-47.

\bibitem{dir2009}
Parlamento Europeo e Consiglio UE (2009).
\textit{Direttiva 2009/138/CE (Solvency~II)}.

\bibitem{smithwilson2001}
Smith, A. e Wilson, T. (2001).
\textit{Fitting Yield Curves with Long-Term Constraints}.
Research Note, Bacon \& Woodrow.

\bibitem{deBoor1978}
de Boor, C. (1978).
\textit{A Practical Guide to Splines}.
Springer-Verlag.

\bibitem{quarteroni2014}
Quarteroni, A., Saleri, F. e Gervasio, P. (2014).
\textit{Calcolo Scientifico}, 6\textsuperscript{a} ed., Springer.

\bibitem{stoer2002}
Stoer, J. e Bulirsch, R. (2002).
\textit{Introduction to Numerical Analysis}, 3\textsuperscript{rd} ed., Springer.

\bibitem{hull2018}
Hull, J. C. (2018).
\textit{Options, Futures, and Other Derivatives}, 10\textsuperscript{th} ed., Pearson.

\end{thebibliography}

\end{document}
